\section{Polynomial images and recurrent switches}
\label{sec:polynomial-volume}

The product schemes of Lemma~\ref{lem:product-schemes} will be
restricted by conditions on the intermediate group elements.  Such restrictions
need not contain an open set, and their images need not be described by
independent parameters.  We therefore need a volume estimate for the
image of an arbitrary measurable parameter set.  We first prove that
estimate, and then show how to retain enough parameter measure through
a fixed finite sequence of switches.

\subsection{Volume of polynomial images}

The scalar sublevel estimate uses a Remez-type argument; see
\cite[Theorems~1.1--1.2]{Ganzburg2017} for classical measurable-set
inequalities and their earlier attribution. The image-volume estimate
below is proved directly, with the coefficient and minor bounds needed
for our varying polynomial maps.

For a set $A\subset\R^b$, write
\[
 |A|_{\mathrm{in}}
   =\sup\{|K|:K\subset A\text{ is compact}\},
\]
where $|\cdot|$ denotes Lebesgue measure in the indicated dimension.
Thus an inner-volume lower bound does not presuppose measurability of
$A$.  All logarithms below are natural.

\begin{lemma}[Polynomial image volume]
\label{lem:polynomial-image}
Fix integers $a\ge b\ge1$ and $D\ge1$.  There are constants
$C>0$ and $A_1,A_2>0$, depending only on $a,b,D$, with the following
property.  Let $\Phi:\R^a\to\R^b$ be a polynomial map of degree at most
$D$, all of whose coefficients have absolute value at most $B\ge0$.
Suppose that one $b$-column minor of $D\Phi$ has a coefficient of absolute
value at least $c>0$.  Put $c_0=\min\{c,1\}$.
For every $R\ge1$ and every Lebesgue-measurable
$E\subset[-R,R]^a$ of measure $m>0$,
\begin{equation}
\label{eq:polynomial-image-quantitative}
 |\Phi(E)|_{\mathrm{in}}
 \ge C c_0^{b+1}(1+B)^{-b^2}
          m^{A_1}(1+R)^{-A_2}.
\end{equation}
One may take, with $d=\max\{1,b(D-1)\}$,
\[
 A_1=1+ad(b+1),\qquad
 A_2=a+a^2d(b+1)+Db^2.
\]

In particular, let $r\to\infty$ along a sequence and let $\Phi_r$
have these fixed dimensions and degree bound.  Suppose their
coefficients converge to those of a polynomial map which is
submersive at some point.  If the $E_r$ are Lebesgue-measurable and
\[
 E_r\subset[-R_r,R_r]^a,\qquad
 1\le R_r\le e^{o(r)},\qquad |E_r|\ge e^{-o(r)},
\]
then
\begin{equation}
\label{eq:polynomial-image-subexponential}
 |\Phi_r(E_r)|_{\mathrm{in}}\ge e^{-o(r)}.
\end{equation}
The same conclusion follows from
\eqref{eq:polynomial-image-quantitative} if the coefficient upper bounds
$B_r$ and a chosen minor-coefficient lower bound $c_r$ satisfy
$\log(1+B_r)+\log(1/\min\{c_r,1\})=o(r)$.
\end{lemma}

\begin{proof}
We first keep a definite part of $E$ away from the zero set of the
chosen derivative minor.  A coordinate slice then retains positive
measure, and a small cube in that slice makes the restricted map
injective.  Change of variables on this cube will give the image-volume
bound.

The first step uses a scalar sublevel estimate.  If a polynomial $Q$ in $k\ge1$
variables has degree at most $d\ge1$ and some coefficient of absolute
value at least $c$, then, for $R\ge1$ and $0<\eta<c$,
\begin{equation}
\label{eq:polynomial-minor-sublevel}
 \big|\{x\in[-R,R]^k:|Q(x)|\le\eta\}\big|
 \le C_{k,d}(1+R)^k(\eta/c)^{1/(kd)}.
\end{equation}
No upper coefficient bound is needed here.

For $k=1$, let the sublevel set have length $l>0$.  It contains
$d+1$ points whose mutual distances are at least $l/[2(d+1)]$:
after fewer than $d+1$ choices, deleting intervals of that radius
about the chosen points removes total length strictly less than $l$.
Lagrange interpolation at these points bounds every coefficient of
$Q$ by
\[
 C_d\eta(1+R)^d l^{-d}.
\]
Indeed the denominators are products of $d$ point differences, while
the coefficients of the numerator polynomials are bounded by
$C_d(1+R)^d$.  The asserted lower bound on one coefficient now proves
\eqref{eq:polynomial-minor-sublevel} in one variable.

For the induction step, write
\[
 Q(x',t)=\sum_{j=0}^d q_j(x')t^j
\]
and choose $q_j$ which contains a coefficient of absolute value at
least $c$.  Put $u=\eta/c\in(0,1)$ and
$\tau=c u^{(k-1)/k}$.  The induction hypothesis bounds the
$(k-1)$-dimensional measure of $\{|q_j|\le\tau\}$ in the base cube
by
\[
 C(1+R)^{k-1}u^{1/(kd)}.
\]
Above its complement, the one-variable estimate bounds the length of
each sublevel slice by
\[
 C(1+R)(\eta/\tau)^{1/d}
   =C(1+R)u^{1/(kd)}.
\]
Fubini's theorem, using the bound $2R$ on the remaining slice
lengths, proves \eqref{eq:polynomial-minor-sublevel}.  The argument
also applies to constant polynomials, since $d$ is only a degree
upper bound.

Return to the minor $Q$ of $D\Phi$ in the statement.  Its degree is at
most $b(D-1)$, so we can use
$d=\max\{1,b(D-1)\}$.  Choose a compact $K\subset E$ with
$|K|\ge m/2$.  For a sufficiently large constant $C_0$, depending
only on $a,b,D$, set
\[
 \eta=c_0\left(\frac{m}{C_0(1+R)^a}\right)^{ad}.
\]
Since $m\le(2R)^a$, we have $0<\eta<c_0\le1$ after increasing
$C_0$.  Estimate \eqref{eq:polynomial-minor-sublevel}, with coefficient
lower bound $c_0$, makes the bad-minor set have measure at most
$m/4$.  Therefore the compact set
\[
 K'=K\cap\{|Q|\ge\eta\}
\]
has measure at least $m/4$.

Permute the domain coordinates so that the chosen minor uses the
first $b$ columns, and write $x=(v,z)$ with $v\in\R^b$.
Fubini supplies a fixed $z$ for which the compact slice $K'_z$
satisfies
\begin{equation}
\label{eq:polynomial-image-slice}
 |K'_z|\ge C_a m(1+R)^{-(a-b)}.
\end{equation}
When $a=b$, there is no slicing and we use $K'$ itself.
For $f(v)=\Phi(v,z)$, its derivative and the Lipschitz constant of its
derivative on $[-R,R]^b$ are both at most
\[
 L=C_{a,b,D}(1+B)(1+R)^D\ge1.
\]

We now retain a small cube on which $f$ is injective.
Cover $[-R,R]^b$ by at most
$C_b((1+R)/\ell)^b$ closed cubes of side at most
\[
 \ell=\frac{\eta}{4\sqrt b\,L^b}.
\]
One such cube $J$ contains at least the reciprocal of this number
times the mass in \eqref{eq:polynomial-image-slice}.  Choose
$v_0\in K'_z\cap J$.  The least singular value of $Df(v_0)$ is at
least $\eta/L^{b-1}$, since its determinant has absolute value at
least $\eta$.  For every $v\in J$,
\[
 \|Df(v)-Df(v_0)\|
 \le L\sqrt b\,\ell
 \le\frac{\eta}{4L^{b-1}}.
\]
Integrating along a line segment in $J$ gives
\[
 |f(v)-f(w)|\ge\frac{3\eta}{4L^{b-1}}|v-w|
       \qquad(v,w\in J).
\]
Thus $f$ is injective on $J$.  Change of variables on its interior,
with the boundary omitted as a null set, gives
\[
 |f(K'_z\cap J)|
   =\int_{K'_z\cap J}|Q(v,z)|\,dv
   \ge\eta |K'_z\cap J|.
\]
The image here is compact and is contained in $\Phi(E)$.  Combining
the cube count and \eqref{eq:polynomial-image-slice} yields
\[
 |\Phi(E)|_{\mathrm{in}}
 \ge C m\eta^{b+1}(1+R)^{-a}L^{-b^2}.
\]
Substituting $\eta$ and $L$ proves
\eqref{eq:polynomial-image-quantitative} with the stated exponents.
This proof uses neither finite global fibers nor a global
multiplicity bound.

Finally, a polynomial map submersive somewhere has a derivative
minor which is a nonzero polynomial.  Coefficient convergence gives
a fixed positive lower bound for one coefficient of that same minor
of $\Phi_r$, for all sufficiently large $r$, and a fixed upper bound
for all coefficients of $\Phi_r$.  Every factor in
\eqref{eq:polynomial-image-quantitative} then has logarithm bounded
below by $-o(r)$ under the hypotheses on $R_r$ and $|E_r|$.
The final variant follows from the same displayed estimate.
\end{proof}

\subsection{Recurrence within a measurable set}

The passage from a large parameter set to a large image is now
available.  The next lemma supplies that parameter set while keeping
every intermediate state in a prescribed measurable set.

For a probability measure $\sigma$ on a group, let
$\check\sigma$ be its image under inversion.  With the convention
used here, a random element of law $\sigma*\check\sigma$ has the
form $s_1s_2^{-1}$ for independent $s_1,s_2$ of law $\sigma$.

\begin{lemma}[Backward recurrence]
\label{lem:backward-recurrence}
Let $J$ be a second countable unimodular locally compact group with
Haar measure $\mu$.  Let $Q\subset J$ be compact with
$0<\mu(Q)<\infty$, and let $S\subset Q$ be Borel with
$\mu(S)>0$.  Fix an integer $h\ge1$ and compactly supported
probability measures $\sigma_1,\ldots,\sigma_h$ on $J$.
Suppose that numbers $D_j\ge1$ satisfy
\[
 \mu(Q\supp\sigma_j)\le D_j\mu(Q)
       \qquad(1\le j\le h).
\]
There are nested Borel sets
\[
 S_0\subset S_1\subset\cdots\subset S_h=S
\]
with positive measure such that, writing
$a_j=\mu(S_j)/\mu(Q)$ and $\tau_j=a_j/(2D_j)$, one has
\begin{equation}
\label{eq:recurrence-mass}
 a_{j-1}\ge\frac{a_j^2}{2D_j}
\end{equation}
and
\begin{equation}
\label{eq:recurrence-transition}
 \mathbb P\{x u_j\in S_j\}\ge\tau_j
       \qquad(x\in S_{j-1}),
\end{equation}
where $u_j$ has law $\sigma_j*\check\sigma_j$.
If the $u_j$ are sampled independently, every $x\in S_0$ therefore
has probability at least $\prod_{j=1}^h\tau_j$ of satisfying
\[
 x u_1\cdots u_j\in S_j\subset S
       \qquad(1\le j\le h).
\]
In a sequence indexed by $r\to\infty$, if $h$ is fixed,
$\mu(S)/\mu(Q)\ge e^{-o(r)}$, and every $D_j\le e^{o(r)}$,
then $\mu(S_0)/\mu(Q)$ and the displayed success probability
are at least $e^{-o(r)}$.
\end{lemma}

\begin{proof}
For a Borel set $A\subset Q$, let
\[
 (T_j1_A)(x)
   =\int\!\int 1_A(x s_1s_2^{-1})
                      \,d\sigma_j(s_1)\,d\sigma_j(s_2).
\]
This is a Borel function.  Right Haar invariance, the substitution
$z=x s_1$, and Fubini give
\begin{align*}
 \int_A T_j1_A(x)\,d\mu(x)
 &=\int_J\left(\int 1_A(zs^{-1})\,d\sigma_j(s)\right)^2d\mu(z)\\
 &\ge\frac{\mu(A)^2}{\mu(Q\supp\sigma_j)}
 \ge\frac{\mu(A)^2}{D_j\mu(Q)}.
\end{align*}
For the first inequality, the inner function is supported in
$Q\supp\sigma_j$, has integral $\mu(A)$, and
Cauchy--Schwarz applies.  Compactness makes the containing set
measurable and of finite measure.

Starting with $S_h=S$, work backwards and define
\[
 S_{j-1}
   =\{x\in S_j:(T_j1_{S_j})(x)\ge\tau_j\},
 \qquad \tau_j=\frac{\mu(S_j)}{2D_j\mu(Q)}.
\]
The contribution from $S_j\setminus S_{j-1}$ to the preceding
integral is at most $\mu(S_j)^2/(2D_j\mu(Q))$.
Since $T_j1_{S_j}\le1$, the remaining contribution implies
\[
 \mu(S_{j-1})\ge\frac{\mu(S_j)^2}{2D_j\mu(Q)},
\]
which proves \eqref{eq:recurrence-mass}; the definition gives
\eqref{eq:recurrence-transition}.  Conditional on the first $j-1$
successful transitions, the current state lies in $S_{j-1}$ and
the next increment retains its prescribed law.  Multiplying the
conditional lower bounds proves the claimed probability estimate.
For fixed $h$, finitely many squarings, products, and divisions by
$e^{o(r)}$ preserve a lower bound of the form $e^{-o(r)}$.
\end{proof}

\subsection{Restricted product parameters}

We combine the two estimates.  Backward recurrence supplies a lower
bound for the probability of successful parameters.  A density bound
converts this to Lebesgue measure, and a submersive limiting polynomial
map converts that measure to endpoint volume.

\begin{corollary}[Restricted product parameters]
\label{cor:restricted-products}
Let $r\to\infty$ along a sequence.  Fix $h\ge1$, positive
integers $a_1,\ldots,a_h$, and $a=\sum_j a_j$.
For each $r$, let $J_r,Q_r,S_r,\mu_r$, and
$\sigma_{j,r}$ satisfy the hypotheses of
Lemma~\ref{lem:backward-recurrence}, with
\[
 \frac{\mu_r(S_r)}{\mu_r(Q_r)}\ge e^{-o(r)},\qquad
 \frac{\mu_r(Q_r\supp\sigma_{j,r})}{\mu_r(Q_r)}
       \le e^{o(r)}.
\]
For each $j$, let $\theta_j\in\R^{a_j}$ have a probability law
$\nu_{j,r}$, and let the $\theta_j$ be independent.  Suppose their
joint law $\nu_r$ gives full measure to a Borel admissible domain
$\Theta_r\subset[-R_r,R_r]^a$, is absolutely
continuous with respect to Lebesgue measure, and satisfies
\[
 1\le R_r\le e^{o(r)},\qquad
 \frac{d\nu_r}{d\theta}\le e^{o(r)}
       \quad\text{almost everywhere}.
\]
Let $v_{j,r}:\R^{a_j}\to J_r$ be Borel maps whose images of the
laws $\nu_{j,r}$ are $\sigma_{j,r}*\check\sigma_{j,r}$.
Finally, let $\Phi_r:\R^a\to\R^b$, $1\le b\le a$, be polynomial
maps of a fixed degree bound whose coefficients converge to a map
submersive somewhere.

There is a starting point $x_r\in S_r$ such that the successful
parameter set
\[
 E_r=\left\{\theta\in\Theta_r:
 x_r v_{1,r}(\theta_1)\cdots v_{j,r}(\theta_j)\in S_r
 \text{ for every }1\le j\le h\right\}
\]
satisfies
\[
 |E_r|\ge e^{-o(r)},\qquad
 |\Phi_r(E_r)|_{\mathrm{in}}\ge e^{-o(r)}.
\]
Here $\Phi_r$ may record a product endpoint, several successive
endpoints, or a tuple of bundle products; it need not record the
intermediate states used in the definition of $E_r$.
\end{corollary}

\begin{proof}
Apply Lemma~\ref{lem:backward-recurrence} and choose
$x_r\in S_{0,r}$.  Independent sampling of the parameters produces
the required independent increment laws, so $\nu_r(E_r)\ge
e^{-o(r)}$.  The set $E_r$ is Borel, because $\Theta_r$, multiplication, and the
increment maps are Borel.  The density upper bound gives
\[
 \nu_r(E_r)\le e^{o(r)}|E_r|,
\]
proving the first conclusion.  Lemma~\ref{lem:polynomial-image}
proves the second.  No independence remains necessary after
restricting to $E_r$.
\end{proof}

\paragraph{Application to nilpotent product schemes.}
In the Lie-group applications, the remaining tasks are to choose output
coordinates that record Haar volume and to verify the limiting rank
condition.  In a simply connected
nilpotent Lie group $M$ of dimension $b$, choose exponential
coordinates so that Haar measure is Lebesgue measure.  If an output
map has the form
\[
 \theta\longmapsto\exp_M(L_r \Phi_r(\theta))
\]
for an invertible linear map $L_r:\R^b\to\Lie(M)$, then the
conclusion becomes an inner Haar-volume lower bound
\[
 |\det L_r|\,e^{-o(r)}.
\]
For a box normalization in $N$, this determinant is bounded above
and below by fixed positive multiples of $e^{rP(E)}$; for a fixed
product of boxes the determinants multiply.  Subsequent left
translations and other Haar-preserving coordinate homeomorphisms leave
this lower bound unchanged.  In particular, one can first record
unconjugated bundle products and only afterwards apply
Haar-preserving conjugations and pass to successive total products.
One does not need bounded normalized coefficients for those later
coordinate changes.

A generator increment with parameter $t$ can be sampled using
$t=s_1-s_2$, where $s_1,s_2$ are independent uniform parameters on
a sufficiently large fixed centered interval.  The resulting
increment law is $\sigma*\check\sigma$.  Replacing every scalar
parameter in a submersive generator product scheme by such a
difference preserves its rank: the linear map
$(s_1,s_2)\mapsto s_1-s_2$ is surjective, and the chain rule
applies.  The parameter intervals can be chosen to contain a
preimage of a submersion point in their interiors.  All parameters
here are the normalized parameters; their number and interval
bounds are fixed before the scale tends to infinity.

\begin{remark}[Scope of the parameter estimate]
\label{rem:restricted-products-scope}
The scheme and its dimensions must be fixed after any needed
subsequence of normalized coordinate laws is chosen, and before
the successful parameter set is restricted.  Each application must
verify a submersive limit for its actual polynomial output map.
Positive successful probability alone is insufficient without
the stated density and rank hypotheses.  For example,
$(x,y)\mapsto(x,0)$ maps a positive-area parameter square to a
zero-area set, while $x\mapsto e^{-r}x$ is submersive for every
$r$ but loses an exponential factor on a fixed interval.

An arbitrary measurable selector of product representations need
not give an absolutely continuous parameter law.  It cannot be
inserted into Corollary~\ref{cor:restricted-products} merely because
its endpoints have positive Haar measure.  When endpoints are
sampled directly by Haar measure, their measure must instead be
retained and used directly in the endpoint packing estimate.
\end{remark}
